\subsubsection{An exact certificate retaining nonzero residuals}
\label{q444:twocross:quantitative-section}

The exclusion has a quantitative extension. It also gives a polynomial
identity before the Gauss or spatial residuals are set to zero. Use the
entire second-frame variables above, with every entry of $v,w$ retained.
More precisely, return here to the matrices and inverse in
\eqref{q444:twocross:eq:second-frame}--\eqref{q444:twocross:eq:original-inverse},
not the literal-frame $H$ used in the preceding determinant comparison.
All of $M,N,m,u,v_0,b,s_1,s_2$ in this subsection are computed in
that same second frame as in \eqref{q444:twocross:eq:s1}--\eqref{q444:twocross:eq:s2}.
Define the full spatial residual matrix and the Hermitian Gauss components
\[
 \mathcal A=s_1v^*+s_2w^*,\qquad
 \gamma_0=et/6+4e(u+2v_0)+8\omega b,
 \qquad \mathcal G=etE-2e(QM+MQ)-4\omega N.
\]
Here $\mathcal A$ is not the original generator $A$: its calligraphic
symbol denotes exactly this two-by-two residual matrix. The full Gauss
residual is $i\operatorname{diag}(\gamma_0,\mathcal G)$, with the
second block of size two. Set $B_2=4b-4\omega e$ and
\[
 \mathcal P=d_M+\frac{e^2}{\omega^2}
       (u+2v_0+t/24+2\omega^2)(v_0+5t/12).
\]
This expression is defined for the original nonzero $\omega$; the
identity below has no such denominator after multiplication.

\begin{proposition}[Polynomial residual certificate and strict error bound]
\label{q444:twocross:quantitative}
For every $v,w\in\mathbb C^2$ and real $e,t,\omega$ in the displayed
second-frame matrices, the following polynomial identity holds:
\begin{align}
 &4\omega^2(\bar m\mathcal A_{22}-v_0\mathcal A_{21})
   +\omega B_2(\bar m\mathcal G_{22}-v_0\mathcal G_{21})\notag\\
 &\hspace{8mm}+e(v_0+5t/12)\bar m\gamma_0\notag\\
 &=4\bar m\bigl[\omega^2d_M+
      e^2(u+2v_0+t/24+2\omega^2)(v_0+5t/12)\bigr].
 \label{q444:twocross:polynomial-certificate}
\end{align}
No scalar, Gauss or spatial equation is assumed in this identity.

Now restrict to the exactly specified set of original tuples with
$Z_1,Z_2\in\mathfrak p$, $t>0$, $\omega\ne0$, for which both
scalar equations and Gauss hold. Let $R_1,R_2\in\mathfrak{su}(3)$
be their two full spatial residuals. Then
\begin{equation}
 \bigl(\|R_1\|_{\mathrm{HS}}^2+
             \|R_2\|_{\mathrm{HS}}^2\bigr)^{1/2}
 \geq
 \frac{\sqrt2\,\mathcal P}
 {\sqrt{\operatorname{tr}M\,(1+v_0^2/|m|^2)}}
 >\frac{5\sqrt2\,e^2t}
 {6\sqrt{\operatorname{tr}M\,(1+v_0^2/|m|^2)}}.
 \label{q444:twocross:residual-lower-bound}
\end{equation}
The full matrix norms are preserved by the two exact unitary frame maps.
In particular, on any such set with $\operatorname{tr}M\leq L$ and
$t\geq t_0>0$, where $L>0$, the strictly positive uniform bound is
\begin{equation}
 \bigl(\|R_1\|_{\mathrm{HS}}^2+
             \|R_2\|_{\mathrm{HS}}^2\bigr)^{1/2}
 >\frac{5\sqrt2\,t_0}
 {3\sqrt{L(1+L^2/\delta^2)}},\qquad \delta=1/(2r).
 \label{q444:twocross:uniform-bound}
\end{equation}
This does not assert that the scalar--Gauss solution set is nonempty.
It proves that no sequence in its displayed bounded portion can make
both spatial residuals tend to zero.
\end{proposition}

\begin{proof}
The complete residual formulas above give
\[
 \mathcal A_{21}=D_2\bar m+B_2N_{21},\qquad
 \mathcal A_{22}=D_2v_0+d_M+B_2N_{22},\quad
 D_2=t/4-\omega^2-4e^2,
\]
and, without any equation imposed,
\[
 \mathcal G_{21}=-6e\bar m-4\omega N_{21},\qquad
 \mathcal G_{22}=-5et/6-8ev_0-4\omega N_{22}.
\]
Subtracting the two indicated spatial entries cancels $D_2$ exactly:
\[
 \bar m\mathcal A_{22}-v_0\mathcal A_{21}
 =\bar m d_M+B_2(\bar mN_{22}-v_0N_{21}).
\]
Multiplication of the indicated Gauss entries gives
\[
 \bar m\mathcal G_{22}-v_0\mathcal G_{21}
 =-e(2v_0+5t/6)\bar m
      -4\omega(\bar mN_{22}-v_0N_{21}).
\]
These cancel the two $N$ entries on the left of
\eqref{q444:twocross:polynomial-certificate}. The remaining factor is
\[
 \gamma_0-2\omega B_2
 =4e(u+2v_0+t/24+2\omega^2).
\]
It yields exactly its stated right side. Each displayed step is a
polynomial identity, including at $m=0$ or $\omega=0$; no inverse or
solution is assigned at those values.

On the stated scalar--Gauss zero set the preceding proof forces the
second-frame form, $e^2=2+\operatorname{tr}M$ and
$\operatorname{Re}m=\delta>0$. Hence $m\ne0$ and
$\operatorname{tr}M>0$. Setting $\gamma_0=\mathcal G=0$ in the
polynomial identity and dividing by its retained nonzero factors gives
\[
 \mathcal A_{22}-\frac{v_0}{\bar m}\mathcal A_{21}
 =\mathcal P>0.
\]
Cauchy--Schwarz on these two entries implies
\[
 \mathcal P\leq
  \sqrt{1+v_0^2/|m|^2}\,\|\mathcal A\|_{\mathrm{HS}}.
\]
Let $S=(s_1\ s_2)$ and $V=(v\ w)$ be the two column matrices.
Then $\mathcal A=SV^*$. Applying Cauchy--Schwarz to each entry and
summing proves $\|SV^*\|_{\mathrm{HS}}^2
\leq\|S\|_{\mathrm{HS}}^2\|V\|_{\mathrm{HS}}^2$.
The right factors are $\|s_1\|^2+\|s_2\|^2$ and
$\operatorname{tr}M$, respectively. Directly from the two off-diagonal
blocks of $X(s)$, $\|X(s)\|_{\mathrm{HS}}^2=2\|s\|^2$.
Thus $R_j=X(s_j)$ proves the first inequality in
\eqref{q444:twocross:residual-lower-bound}.

The exact expansion
\begin{align*}
 &(u+2v_0+t/24+2\omega^2)(v_0+5t/12)\\
 &=5t\omega^2/6
   +v_0(u+2v_0+t/24+2\omega^2)
   +(5t/12)(u+2v_0+t/24)
\end{align*}
has nonnegative remainder terms, the last strictly positive for $t>0$.
Together with $d_M\geq0$ and $e^2>0$ it gives
$\mathcal P>5e^2t/6$, proving the strict second inequality.
Finally $e^2=2+\operatorname{tr}M\geq2$,
$0\leq v_0\leq\operatorname{tr}M\leq L$, and $|m|\geq\delta$
give \eqref{q444:twocross:uniform-bound}. Since unitary conjugation
preserves the Hilbert--Schmidt norm, the same inequalities hold for
the complete spatial residuals in the original literal matter frame.
\end{proof}
